\begin{lemma}[Odd order from an elementary closed missing color]
Let $G=(V,E)$ be a finite loopless multigraph, let $e_0\in E$, let
$\varphi$ be a proper partial coloring with colors in $[k]$, and let $T$
be a tree-sequence rooted at $e_0$.  Suppose that $V(T)$ is elementary
with respect to $\varphi$.  If
$\alpha\in\overline\varphi(T)$ and $T$ is
$\{\alpha\}$-closed under $\varphi$, then $|V(T)|$ is odd.
\end{lemma}
\begin{proof}
Write $X=V(T)$.  By \noderef{TreeMissingColors}, choose
$x\in X$ at which $\alpha$ is missing.  Elementarity in the sense of
\noderef{ColoredVertexSetElementary} implies that $\alpha$ is not missing
at any other vertex of $X$: otherwise the two missing-color sets would
intersect in $\alpha$.  Hence, by \noderef{PresentColors} and
\noderef{MissingColors}, every vertex of $X-\{x\}$ is incident with an
exact edge identity colored $\alpha$, whereas $x$ is incident with none.

For every $v\in X-\{x\}$, that incident $\alpha$-edge is unique.  Indeed,
two distinct such identities would be adjacent at $v$ and would receive
the same color, contrary to properness as defined in
\noderef{IsProperEdgeColoring}.  Moreover, every $\alpha$-edge incident
with a vertex of $X$ has its other endpoint in $X$.  Otherwise it would
be an exact member of the $\alpha$-colored boundary set from
\noderef{BoundaryColorEdges}, contradicting the assumed
$\{\alpha\}$-closedness in \noderef{TreeColorSetClosed}.

Let $M$ be the set of exact $\alpha$-colored edge identities whose two
endpoints lie in $X$.  The preceding paragraph shows that the endpoint
sets of the edges in $M$ cover exactly $X-\{x\}$.  These endpoint sets are
pairwise disjoint, because two distinct members of $M$ sharing an endpoint
would again violate properness.  Each member has exactly two distinct
endpoints by looplessness in \noderef{LooplessMultigraph}.  Therefore
\[
 |X|-1=|X-\{x\}|=2|M|,
\]
so $|X|-1$ is even and $|X|=|V(T)|$ is odd.
\end{proof}
